% FRL letter version — companion to the 40-page working paper (online appendix).
% Every number in this letter is drawn from the archived, test-verified result files
% (scripts/test_numbers.py, 290 assertions; scripts/frl_letter_check.py for this letter).
\documentclass[11pt]{article}
\usepackage[margin=1in]{geometry}
\usepackage{newtxtext,newtxmath}
\usepackage{booktabs}
\usepackage{setspace}
\usepackage[hidelinks]{hyperref}
\usepackage{caption}
\usepackage{amsmath}
\usepackage{microtype}
\usepackage{natbib}
\captionsetup{font=small,labelfont=bf,skip=4pt}
\setlength{\bibsep}{3pt}
\setlength{\parindent}{0pt}
\setlength{\parskip}{5pt}
\renewcommand{\baselinestretch}{1.05}

\title{\bfseries Dissecting LLM News Sentiment Scores: De-Identification, Batch Composition, Scoring Noise, and Price Echoes in Return Predictability\thanks{A companion working paper (40 pp.) with all appendices, the mechanism accounting, power tables, preregistered design documents, and a full replication package is available as online appendix at \url{https://github.com/43aquarius/dissecting-llm-news-sentiment-scores}. This letter is a compressed version of that working paper; all references to sections, tables, and appendices refer to it.}}
\author{Chang Tan\thanks{Dalian University of Technology. Email: aquars43@foxmail.com.}}
\date{October 2026}

\begin{document}
\maketitle

\begin{abstract}
\noindent Search-collected news headlines on 36 U.S.\ large-cap firms (341 in-window, July--October 2026), scored $-100$ to $+100$ by GLM-4-Plus across six sittings and fourteen re-scorings, deliver a marginal and draw-dependent next-day association ($t=2.10$, BH $q=0.062$)---and a measurement audit with robust findings. De-identification, replacing the firm's name with a placeholder, raises score levels by two to three points; swapping in another firm's name moves levels by essentially nothing: the effect is the removal of the name, not a company prior. Batch composition, which headlines share an API call, moves levels by 3.6 points, survives inverted run order, and defies a linear-anchoring mechanism; it leaves slopes intact but doubles the trade count of a fixed $+50$ threshold. The instrument is stochastic (ICC $0.92$; 30\% exact replication across four calls; temperature zero restores bitwise determinism for only 83\% of headlines). Predictability directionally concentrates in price-echo headlines ($p=0.10$), and news coverage is itself selected on realized volatility ($q=0.014$). Four identity manipulations leave predictive content formally equivalent under two one-sided tests. A seven-step audit protocol should precede any deployment of LLM sentiment as a predictor. \\[2pt]
\textbf{JEL:} G10; G12; C52; C55. \textbf{Keywords:} large language models; sentiment analysis; return predictability; measurement validity; look-ahead bias.
\end{abstract}

\section{Introduction}

LLM-based sentiment scoring of firm news has moved from novelty to infrastructure \citep{lopezlira2026,kirtac2024}. The worry list is growing alongside it: look-ahead bias in GPT sentiment signals \citep{glasserman2023}, memorization of realized returns \citep{lopezlira2025,didisheim2025}, miscalibrated extrapolation from historical prices \citep{chen2025}, sensitivity to rating-scale formats \citep{jo2026}, and a call to treat bias evaluation as an explicit requirement for finance LLMs \citep{kong2026}. What is still missing is a within-design, joint quantification of the \emph{non-semantic} channels that a sentiment score rides on: whose name is in the text, which headlines share the API call, and how much of the score is sampling noise. This letter provides that audit on one design and one instrument, and reports what survives.

We audit a short-horizon predictability design in the spirit of \citet{lopezlira2026}: 542 search-collected headlines about 36 U.S.\ large-cap firms (341 dated inside July 1--October 14, 2026), scored on a $-100$ to $+100$ scale by GLM-4-Plus under conditions that each vary one ingredient---identity visible, removed, or swapped; repeat calls; controlled and randomized batch composition; temperature-zero decoding; and the Loughran--McDonald dictionary \citep{loughran2011} as an external benchmark. Six sittings and fourteen re-scorings make every headline-level channel reproducible with call-clustered inference. The conceptual frame is that the score is not a thermometer reading: it mixes text semantics with what the scorer brings to the text---identity references, the batch context, sampling noise---and what the sampling frame brought to the text---retrospective price language and coverage selection. Each ingredient gets its own contrast, its own clustered inference, and its own equivalence assessment where the claim is invariance.

The baseline next-day association is marginal and draw-dependent: the pre-specified coefficient is $0.0149$ ($t=2.10$; wild-cluster bootstrap $p=0.030$; Benjamini--Hochberg $q=0.062$ over the declared seven-comparison family), $t=1.82$ with a same-day return control, and $t\in[1.47,2.12]$ across sixteen archived named draws. The design's resolution is bounded: at 80\% power the minimum detectable coefficient is $0.020$, a third larger than the estimate itself, and the power against the observed effect is $0.56$. We do not claim a new anomaly; the honest headline is the measurement.

Four findings are robust. First, de-identification raises score levels by two to three points, and the name-swap control indicates the effect is the \emph{removal} of the named reference rather than a material company-specific prior. Second, batch composition moves levels by $+3.6$ points, reproduces with the run order inverted, and is not explained by linear anchoring; it leaves regression slopes and cross-sectional ranks intact but doubles the trade count of a fixed threshold and flips the sign of day-over-day changes in daily sentiment on up to a quarter of trading days. Third, the instrument is stochastic: ICC $0.92$ across sessions, 30\% exact replication across four calls, and temperature zero does not restore determinism. Fourth, under a formal equivalence battery, predictive content is unchanged by every identity manipulation, the LLM-versus-dictionary contrast is bidirectionally undecidable at this sample size, and the echo-free subsample cannot be declared equivalent to zero. To our knowledge, the identity, batch-composition, and sampling-noise channels have not previously been quantified jointly in the LLM-finance setting---NLP work shows that scores move with batch membership, item position, and option context in non-finance tasks \citep{lin2024batch,liu2024lost,zheng2024,korikov2025}---and the bias taxonomy of \citet{kong2026} names them as threats to financial-LLM evaluation; we measure them.

\section{Data and audit design}

The corpus is 542 search-collected headlines on 36 U.S.\ large-caps; 341 are dated in-window, and a relevance filter leaves 335. Firm-day aggregation yields 240 news firm-days (36 firms, 59 trading days); the regression sample is 205 firm-day observations with non-missing lags. The coverage-selection test uses the full 2,664-firm-day panel (2,628 with same-day returns). Each headline is scored by GLM-4-Plus (API, default decoding, batches of 25, so one pass is 14 correlated API calls) returning a sentiment score and a relevance rating; per-call logs record the served model. Two design details matter for what follows. First, sittings are hours to days apart, and a naive cross-sitting contrast therefore mixes the treatment with serving-stack drift; the interleaved and rotated sessions exist to strip that confound out. Second, anonymization is a procedure, not a promise: a verification pass caught ten headlines whose anonymized text still named the focal firm (five of them \emph{Visa Bulletin} immigration notices, a corpus quirk we disclose and bound by exclusion), and all headline-level conclusions are reported with and without the affected corpus.

Five sittings field the audit. Session 1 (Oct 15) scores the named corpus. Session 2 (Oct 16) adds an anonymized pass and three named repeats. Session 3 (Oct 18) interleaves named and anonymized rounds within batches---alternating condition order across batches---and adds a composition control that re-scores the identical 341 headlines inside the original 542-article batches. Session 4 inverts the control's run order (composition pass first) and adds a seeded random re-batch of the same corpus. Session 5 (Oct 19) fields six conditions in rotated batch order: named, corrected anonymization (verified to leave zero residual firm-name mentions), name-swap (each firm's name replaced by another panel firm's), a neutral prefix, and two temperature-zero passes.

Inference is pre-specified and layered: two-way clustered regressions; a day-level wild-cluster bootstrap; day-clustered bootstrap difference tests; and call-batch-clustered $t$-statistics (14--22 clusters, with round-level checks) for level contrasts, because the 25 scores inside one API response share a response-level mood that headline-level errors would ignore. The declared five-hypothesis, seven-comparison family is Benjamini--Hochberg corrected; every staged addition (interleaved session, mechanism, session 5, equivalence) is reported outside the family with its design document archived before data collection, and the staged design's forking-path exposure is disclosed. Equivalence is assessed by percentile two one-sided tests on day-clustered bootstrap 90\% intervals against a margin of half the benchmark coefficient, with a margin-free endpoint statistic also reported; this battery is exploratory and labeled as such. Loughran--McDonald dictionary scores on the same titles (nonzero on only 16.2\% of firm-days) provide the external benchmark. All scores, logs, seeds, and scripts are archived; the replication package regenerates every number in this letter.

\section{Findings}

Table~\ref{tab:audit} summarizes the audit; Table~\ref{tab:tost} the equivalence battery.

\paragraph{De-identification shifts levels; the name-swap shows it is not a prior.}
In the interleaved session, anonymized scores sit $+2.8$ points above named scores on the identical headlines (call-batch-clustered $t=2.73$, $p=0.017$; round-level $t=6.82$; all four round pairs negative). Session 5's corrected anonymization recurs at $+2.1$ (Wilcoxon $p<0.001$; batch-clustered $p=0.175$ with 14 clusters)---direction and magnitude consistent with the interleaved estimate. The discriminating control is the name-swap: replacing company A's name with company B's moves levels by $+0.3$ ($t=0.40$), while the placeholder differs from the swapped-name condition by $2.4$ points ($p=0.001$; batch-clustered $p=0.08$). The reading is that a specific name is removed---a placeholder-rewording effect---rather than which company the name denotes. A pre-registered sixth session adds the missing cell---invented firm names, natural text, no real-world prior---and lands almost exactly midway between the named and anonymized levels ($+1.20$ above named, $-1.33$ below anonymized; neither clustered-significant, $n=341$), so the displacement splits about evenly between a real-name component and a placeholder-rewrite component, while the de-identification effect itself replicates a third time ($-2.54$, clustered $p=0.15$). A level-layer equivalence battery at a pre-fixed 1.25-point margin certifies only the prefix control ($-0.13$ in point estimate, minimum passing margin 1.10 points); the swap's minimum passing margin is 1.73 points headline-level---small, 13\% of the placeholder effect in point estimate, but not certifiable at half the effect at this sample size. A leakage audit of the swapped texts finds 15.0\% retaining an original-firm lexical cue (a broader rule: 26.7\%) and six headlines (1.8\%) retaining the identity outright through concatenated name variants the verification gate cannot see; excluding every flagged headline leaves the swap contrast at $+0.36$ to $+0.42$---the null is not an artifact of identity leakage. The naive two-session benchmark overstates the effect twice over. First, it confounds identity with a $+5.79$-point cross-session drift on the unchanged named protocol ($t=7.03$) that the interleaved design decomposes into $+1.05$ (endpoint, 542-composition) $+3.61$ (composition, within session) $+1.14$ (endpoint, 341-composition), with the pure endpoint component between sessions 3 and 4 at only $+0.56$ ($t=0.31$). Second, its $+4.75$ level difference was measured before the anonymization verification pass existed; the corrected, verified contrast is the $+2.1$ to $+2.8$ reported above.

\paragraph{Batch composition is a level channel with boundary costs.}
Re-scoring the same 341 headlines inside the original 542-article batches raises levels by $+3.61$ points in session 3 (clustered $t=2.74$, $p=0.012$), by $+3.50$ with the run order inverted in session 4 (clustered $t=2.17$), pooled $+3.56$ ($z=3.4$); re-batching the same corpus alone moves levels by $+4.72$ (clustered $t=2.02$, $p=0.065$---suggestive). The mechanism is bounded, not closed. Linear anchoring on batch-mates' fixed content is estimated at $\theta=+0.36$ ($t=4.9$), a telescoping identity proves it cannot shift the mean of a fixed corpus under re-partitioning, and for the composition contrast the linear part predicts $+0.04$ points against the $+3.56$ observed---about 1\%---because the added out-of-window titles average $-1.3$ points \emph{below} the headlines they sit beside, the wrong direction for a full explanation. The estimated non-linearity does not close the residual either: positive batch-mates assimilate at $+0.69$ ($t=2.3$) while negative ones do nothing, which would predict a small \emph{decrease} here because the added titles are less positive than those they replace. Downstream, the level shift leaves the predictive slope inside $[0.010,0.016]$ across seven archived sources (every composition or re-batch contrast of slopes has $p>0.9$) and cross-sectional ranks correlate at $\rho\geq0.81$; but a fixed ``score $\geq+50$'' rule more than doubles its trade count (22 to 50 firm-days in session 3; 27 to 43 in session 4), daily means shift by 4.0--4.6 points, and day-over-day sign changes flip on 16--24\% of trading days. Composition is harmless to conditional, cross-sectional uses and disqualifying for threshold and time-series uses. A pre-registered dose-response pass (the same 50 headlines scored alone, in fives, and in twenty-fives within one sitting) locates the zero point: the singleton level matches the 25-batch level on identical titles to $-0.2$ points ($t=-0.09$), so the composition channel requires a changed mate pool, not batch context in general.

\paragraph{The instrument is stochastic, and temperature zero does not close it.}
Four calls on identical headlines give $\mathrm{ICC}(2,4)=0.919$ pooled across sessions ($0.949$ within), 29.9\% exact replication, and a within-headline standard deviation of 7.8 points; 13\% of residual variance is a response-level mood shared by the 25 scores of one API response, which is why level contrasts need call-clustered inference. Two temperature-zero passes return bitwise-identical scores for 83.3\% of headlines (284 of 341; within-headline SD 4.8 versus 7.2 under default decoding) and sit $+1.06$ points above the default level ($p=0.052$): a pattern consistent with server-side nondeterminism. A single-call pipeline embeds a $\pm8$-point noise floor per headline into every downstream construct \citep{herrera2025}.

\paragraph{Price echoes and coverage selection bias the return layer.}
15\% of firm-days carry headlines echoing realized prices; their next-day coefficient is $0.047$ ($t=1.80$) against $0.003$ ($t=0.61$) for echo-free firm-days, a difference the day-clustered bootstrap prices at $p=0.10$ with power $0.31$---directional evidence, reported as such. A rulebook audit of the 52 flagged headlines finds 82.7\% precision, but it is LLM-assisted, not independent human annotation, and the failures are instructive: three matches are retail promotions and three more are revenue rather than price language. Coverage itself is selected: news firm-days exhibit mean absolute same-day returns of 215~bps against 154~bps, a gap a pooled two-sample test would call overwhelming ($t=3.12$) but which requires panel inference---the same firms recur across days and all firms share each day's market shock---and a firm-and-day fixed-effects test on the full panel gives $t=2.76$ ($p=0.006$; BH $q=0.014$; day-cluster bootstrap CI $[25,103]$~bps). The audit sample over-weights high-activity events; this is detection, not correction.

\paragraph{Predictive content is formally equivalent under all identity manipulations.}
Four de-identification contrasts pass two one-sided tests at a margin of half the benchmark coefficient, with 90\% interval endpoints at 21--42\% of the benchmark and power 0.74--0.99 (Table~\ref{tab:tost}); the name-swap contrast is the edge case (56\%, power 0.75). The LLM-versus-dictionary contrast is bidirectionally undecidable: its endpoint statistic reaches 4.4 times the benchmark and 80\% power would require roughly 165 times the sample---with the dictionary nonzero on only 16.2\% of firm-days, this comparison cannot settle whether the LLM adds information over the dictionary, in either direction. The priced-versus-clean contrast would require 144 times. And the echo-free coefficient, while insignificant, is not established as equivalent to zero (endpoint at 91\% of the benchmark): the design can bound the price-echo channel only directionally at this sample size. Non-rejection is not equivalence; where we claim invariance, we claim it through equivalence tests. The verdicts are also margin-specific: at a stricter economics-anchored margin of 10 bps of next-day response per one-SD sentiment shock---independent of the benchmark $t$, which ranges over $[1.47,2.12]$ across draws---no contrast passes (minimum margins 11.5--29.8 bps/SD across the identity contrasts), so invariance is certified at roughly a fifth to half of the benchmark effect, and no tighter.

\begin{table}[t]
\centering\small
\caption{The measurement audit at a glance (score points unless noted; call-batch-clustered inference for level contrasts)}
\label{tab:audit}
\begin{tabular}{llll}
\toprule
Channel & Contrast & Estimate & Inference \\
\midrule
De-identification & named vs.\ placeholder, interleaved & $+2.8$ & $t=2.73$, $p=0.017$; rounds $t=6.82$ \\
Name-swap & firm A's name vs.\ firm B's & $+0.3$ & $t=0.40$, n.s. \\
Batch composition & 542- vs.\ 341-batch, same session & $+3.56$ pooled & $z=3.4$; inverted order $t=2.17$ \\
Re-batch alone & same corpus, re-partitioned & $+4.72$ & $t=2.02$, $p=0.065$ \\
Repeat-call noise & four calls on identical text & ICC $0.92$--$0.95$ & 29.9\% exact; within-SD 7.8 \\
Temperature zero & two identical passes & 83.3\% bitwise & 57 of 341 move; SD 4.8 vs.\ 7.2 \\
Price echo & priced vs.\ clean next-day $\beta$ & $0.047$ vs.\ $0.003$ & difference $p=0.10$, power 0.31 \\
Coverage selection & covered vs.\ uncovered $|$ret$|$ & 215 vs.\ 154 bps & FE $t=2.76$, $q=0.014$ \\
\bottomrule
\end{tabular}
\end{table}

\begin{table}[t]
\centering\small
\caption{Equivalence battery (percentile TOST, day-clustered bootstrap 90\% CIs, $B=9{,}999$; margin $\delta^{*}=0.0074=$ half the benchmark named coefficient $0.0149$; exploratory, outside the declared family). EMT is the 90\% interval endpoint for the coefficient difference, as a percentage of the benchmark. ``Undecidable'' marks contrasts that cannot pass or fail at this sample size.}
\label{tab:tost}
\begin{tabular}{lccc}
\toprule
Contrast & EMT & Power at $\delta^{*}$ & Verdict at $\delta^{*}$ \\
\midrule
Named vs.\ anonymized (cross-session) & 27\% & 0.90 & equivalent \\
Named vs.\ anonymized (same-session) & 21\% & 0.99 & equivalent \\
Named vs.\ anonymized (interleaved) & 42\% & 0.74 & equivalent \\
Session-5 named vs.\ corrected anonymized & 35\% & 0.91 & equivalent \\
Session-5 named vs.\ name-swapped & 56\% & 0.75 & not established \\
Named vs.\ Loughran--McDonald dictionary & 442\% & 0.00 & undecidable ($\approx$165$\times$ sample) \\
Priced vs.\ clean & 655\% & 0.00 & undecidable ($\approx$144$\times$ sample) \\
Clean subsample vs.\ zero & 91\% & 0.00 & not established ($\approx$6$\times$ sample) \\
\bottomrule
\end{tabular}
\end{table}

\section{Protocol, limitations, and conclusion}

The audit protocol that falls out is seven lines: anonymize the text and verify the anonymization; field identity contrasts interleaved in one sitting; hold batch composition fixed across any comparison; score at least twice and report an ICC or repeat-call dispersion; report how headline results depend on the scoring draw; flag retrospective price language and test coverage selection on the panel; and claim invariance only through equivalence tests, disclosing model, vendor, and session structure. Each step is cheap---here, the entire audit used $\sim$3{,}400 archived API calls on 542 headlines, the echo flag is a five-line regular expression, and the selection test is a two-line fixed-effects regression---and each one changed a conclusion in this project, including catching a ten-headline anonymization failure and a session-drift masquerade inside our own first draft \citep{christensen2026,bisharat2026,benhenda2026}.

The limitations are structural. All findings come from one model through one vendor SDK (per-call logs and probe tests confirm every request was routed to GLM-4-Plus), so this is a within-instrument case study; LLM sentiment instruments themselves differ across models \citep{teles2025}, a low-cost cross-vendor replication design is in the appendix, and we do not extrapolate to other models. The core regression has 205 observations; headline-level point timestamps are unavailable, so the return-window alignment is daily; the price-echo audit is LLM-assisted rather than independently human-annotated; and the staged design, while pre-documented at each stage, accumulates forking-path exposure that we disclose rather than claim away \citep{romano2005}.

LLM news-sentiment scores contain measurable, replicable components that have nothing to do with what the text means: a two-to-three-point de-identification effect that is not a company prior, a 3.6-point batch-composition effect that survives order inversion and defies the linear mechanism, and an eight-point sampling floor that temperature zero does not close. None of it moves the next-day predictive coefficient---formally. All of it moves levels, thresholds, and time-series signals. Audit the instrument before deploying the signal.

\begin{spacing}{0.95}
\small
\begin{thebibliography}{20}
\setlength{\itemsep}{1pt}

\bibitem[Benhenda(2026)]{benhenda2026}
Benhenda, M., 2026, Look-ahead-bench: A standardized benchmark of look-ahead bias in point-in-time LLMs for finance, working paper, arXiv:2601.13770.

\bibitem[Bisharat and Hean(2026)]{bisharat2026}
Bisharat, A., and O. Hean, 2026, Evaluating financial sentiment in the age of AI, working paper, SSRN 7177498.

\bibitem[Chen et~al.(2024)]{chen2025}
Chen, S., T.~C. Green, H. Gulen, and D. Zhou, 2024, What does ChatGPT make of historical stock returns? Extrapolation and miscalibration in LLM stock return forecasts, working paper, arXiv:2409.11540.

\bibitem[Christensen and Hansen(2026)]{christensen2026}
Christensen, T., and S. Hansen, 2026, Performing valid inference with AI/ML-generated covariates: A guide for empirical practice, \emph{AEA Papers and Proceedings} 116, 92--97.

\bibitem[Didisheim et~al.(2025)]{didisheim2025}
Didisheim, A., M.~G. Fraschini, and L. Somoza, 2025, AI's predictable memory in financial analysis, \emph{Economics Letters} 256, 112602.

\bibitem[Glasserman and Lin(2024)]{glasserman2023}
Glasserman, P., and C. Lin, 2024, Assessing look-ahead bias in stock return predictions generated by GPT sentiment analysis, \emph{Journal of Financial Data Science} 6(1), 25--42.

\bibitem[Herrera-Poyatos et~al.(2025)]{herrera2025}
Herrera-Poyatos, D., C. Pel\'aez-Gonz\'alez, C. Zuheros, A. Herrera-Poyatos, V. Tejador, F. Herrera, and R. Montes, 2025, An overview of model uncertainty and variability in LLM-based sentiment analysis, \emph{Frontiers in Artificial Intelligence} 8, 1609097.

\bibitem[Jo et~al.(2026)]{jo2026}
Jo, W., G. Kim, and M.~S. Hee, 2026, Change the number, change the answer: Assessing rating scale sensitivity in LLM-based text scoring, working paper, arXiv:2610.22187.

\bibitem[Kirtac and Germano(2024)]{kirtac2024}
Kirtac, K., and G. Germano, 2024, Sentiment trading with large language models, \emph{Finance Research Letters} 62, 105227.

\bibitem[Kong et~al.(2026)]{kong2026}
Kong, Y., H. Lee, Y. Hwang, A. Lopez-Lira, B. Levy, D. Mehta, Q. Wen, C. Choi, Y. Lee, and S. Zohren, 2026, Evaluating LLMs in finance requires explicit bias consideration, working paper, arXiv:2602.14233.

\bibitem[Korikov et~al.(2025)]{korikov2025}
Korikov, A., P. Du, S. Sanner, and N. Rekabsaz, 2025, Batched self-consistency improves LLM relevance assessment and ranking, \emph{Proceedings of EMNLP}, arXiv:2505.12570.

\bibitem[Lin et~al.(2024)]{lin2024batch}
Lin, J., M. Diesendruck, L. Du, and R. Abraham, 2024, BatchPrompt: Accomplish more with less, \emph{Proceedings of ICLR}, arXiv:2309.00384.

\bibitem[Liu et~al.(2024)]{liu2024lost}
Liu, N.~F., K. Lin, J. Hewitt, A. Paranjape, M. Bevilacqua, F. Petroni, and P. Liang, 2024, Lost in the middle: How language models use long contexts, \emph{Transactions of the Association for Computational Linguistics} 12, 157--173.

\bibitem[Lopez-Lira et~al.(2025)]{lopezlira2025}
Lopez-Lira, A., Y. Tang, and M. Zhu, 2025, The memorization problem: Can we trust LLMs' economic forecasts?, working paper, arXiv:2504.14765.

\bibitem[Lopez-Lira and Tang(2026)]{lopezlira2026}
Lopez-Lira, A., and Y. Tang, 2026, Can ChatGPT forecast stock price movements? Return predictability and large language models, \emph{Journal of Financial Economics} 184, 104335.

\bibitem[Loughran and McDonald(2011)]{loughran2011}
Loughran, T., and B. McDonald, 2011, When is a liability not a liability? Textual analysis, dictionaries, and 10-Ks, \emph{Journal of Finance} 66, 35--65.

\bibitem[Romano and Wolf(2005)]{romano2005}
Romano, J.~P., and M. Wolf, 2005, Stepwise multiple testing as formalized data snooping, \emph{Econometrica} 73, 1237--1282.

\bibitem[Teles and Figueiredo(2025)]{teles2025}
Teles, L.~E.~P., and C.~M.~S. Figueiredo, 2025, Comparing LLMs for sentiment analysis in financial market news, working paper, arXiv:2510.15929.

\bibitem[Zheng et~al.(2024)]{zheng2024}
Zheng, C., H. Zhou, F. Meng, J. Zhou, and M. Huang, 2024, Large language models are not robust multiple choice selectors, \emph{Proceedings of ICLR}, arXiv:2309.03882.

\end{thebibliography}
\end{spacing}

\end{document}
